% Cinemática inversa no plano do braço: ponto do pulso e braço de 2 DOF
% Esboço (pose ilustrativa); unidades em cm; origem no ombro (J2)
\begin{tikzpicture}[scale=0.32, >=stealth,
    elo/.style={line width=2.2pt, line cap=round},
    junta/.style={circle, draw, fill=white, line width=1pt, inner sep=2.2pt},
    cota/.style={gray!70!black, thin},
    ang/.style={->, thin}]
  % pontos: ombro (J2), cotovelo (J3), pulso (J4), J5 e alvo (centro de preensão)
  \coordinate (J2) at (0,0);
  \coordinate (J3) at (3.8,10.3);
  \coordinate (J4) at (13.2,8.7);
  \coordinate (J5) at (20.5,8.7);
  \coordinate (G)  at (20.5,-6.4);

  % eixos com origem no ombro
  \draw[->, thin] (-3,0) -- (25,0) node[below left] {$d$};
  \draw[->, thin] (0,-8) -- (0,15) node[left] {$h - L_1$};

  % último elo (imposto pelo ângulo do pulso), a tracejado
  \draw[line width=1.6pt, dashed, gray] (J4) -- node[above] {$L_4$} (J5) -- node[right] {$D$} (G);
  \fill (G) circle (6pt); \node[below right=1pt and 3pt] at (G) {alvo $(d,\, z - L_1)$};
  \draw[cota, dashed] (J5) -- ++(4.5,0);
  \draw[ang] ($(J5)+(3,0)$) arc[start angle=0, end angle=-90, radius=3] node[pos=0.5, right=2pt] {$\alpha_3$};

  % braço de 2 DOF
  \draw[elo] (J2) -- node[left] {$L_A = L_2$} (J3);
  \draw[elo] (J3) -- node[above] {$L_B = L_3$} (J4);
  \draw[cota, dashed] (J2) -- node[pos=0.6, below right=-1pt, black] {$r_3$} (J4);

  % ângulos
  \draw[ang] (6,0) arc[start angle=0, end angle=33.4, radius=6] node[pos=0.5, right=1pt] {$\alpha_1$};
  \draw[ang] (33.4:4) arc[start angle=33.4, end angle=70, radius=4] node[pos=0.5, above right=0pt] {$\beta_1$};
  \draw[ang] ($(J3)+(249.8:2.5)$) arc[start angle=249.8, end angle=350.3, radius=2.5] node[pos=0.55, below=1pt] {$\beta_2$};

  % juntas
  \node[junta] at (J2) {}; \node[below left=2pt] at (J2) {J2 (ombro)};
  \node[junta] at (J3) {}; \node[above=4pt] at (J3) {J3 (cotovelo)};
  \node[junta] at (J4) {}; \node[above=5pt] at (J4) {J4};
  \node[below right=2pt and -2pt, font=\footnotesize] at (J4) {$(x_{3d},\, y_{3d})$};
  \node[junta] at (J5) {}; \node[above right=3pt and 0pt] at (J5) {J5};
\end{tikzpicture}